% !TEX root = ../thesis.tex

\chapter{Introduction} \label{chp:intro}

\section{Background and Motivation}

Legged robots have demonstrated remarkable progress in locomotion over the past decade,
achieving robust walking, running, and stair-climbing on real-world terrain through
model-based control and, more recently, reinforcement learning (RL)~\cite{lee2020learning,kumar2021rma}.
However, locomotion alone severely limits the robot's practical utility.
Many real-world deployment scenarios---warehouse logistics, search-and-rescue, and
domestic assistance---require a robot to not merely navigate its environment but also
to interact with objects within it.
Attaching a robotic arm to a legged platform gives rise to the problem of
\emph{legged loco-manipulation}: the robot must simultaneously maintain stable
locomotion and execute purposeful manipulation, two tasks that impose competing
demands on the body posture, base dynamics, and joint torques.

A fundamental prerequisite for effective arm-based grasping on a mobile legged
platform is a stable, compliant body base.
The robotic arm is mounted on the torso of the quadruped; any unexpected roll,
pitch, or yaw disturbance of the body directly propagates through the mechanical
linkage into end-effector positioning error.
Consequently, before the arm can reliably reach a target, the locomotion policy
must (i)~maintain a near-level body orientation, (ii)~reject perturbations rapidly,
and (iii)~respond to body-pose commands issued by the arm policy so that the
entire body can lean into difficult reach envelopes.
This tight coupling between locomotion and manipulation motivates a cooperative,
multi-policy control architecture in which a \emph{locomotion policy} and an
\emph{arm policy} are trained to communicate and co-adapt.

This project reproduces and extends the \emph{RoboDuet} framework proposed by
Pan~et~al.~\cite{pan2024roboduet}, which introduced exactly such a cooperative
whole-body controller for legged loco-manipulation.
The original work demonstrated the approach on the Unitree Go1 quadruped body
paired with the Unitree Z1 arm.
Our contribution is the systematic adaptation of this framework to the heavier
Unitree~\textbf{B2} quadruped equipped with the same Z1 arm, trained entirely
in the NVIDIA Isaac Gym~\cite{makoviychuk2021isaac} physics simulator using
Proximal Policy Optimisation~(PPO)~\cite{schulman2017proximal}.
In doing so, we identify and resolve the key engineering challenges that arise
when porting a locomotion-oriented RL framework from a lighter platform (Go1,
$\approx\!12$\,kg) to a significantly heavier one (B2, $\approx\!52$\,kg).

\section{Problem Statement}

Given the B2+Z1 whole-body system, the goal is to train a control policy that
enables the quadruped to walk on flat terrain while simultaneously tracking
6-DoF end-effector pose targets with the Z1 arm.
More precisely, the desired behaviour is:
\begin{enumerate}
  \item \textbf{Stable locomotion:} The quadruped follows commanded forward
        velocity, lateral velocity, and yaw-rate targets without falling or
        exhibiting excessive oscillation.
  \item \textbf{End-effector pose tracking:} The Z1 arm tracks randomly sampled
        Cartesian positions (reach length $l$, polar angle $p$, azimuthal angle $y$)
        and end-effector orientations (roll, pitch, yaw in the end-effector frame).
  \item \textbf{Whole-body coordination:} The arm policy communicates desired body
        roll and pitch adjustments to the locomotion policy, allowing the torso to
        lean in order to extend the arm's reachable workspace.
\end{enumerate}

\section{Contributions}

The main contributions of this work are as follows:
\begin{enumerate}
  \item We identify and correct a critical PD-gain mismatch that prevented stable
        B2 locomotion when directly transplanting the Go1-tuned framework
        (\S\ref{sec:framework_adaptation}).
  \item We implement and validate a three-stage, curriculum-based RL training
        pipeline specifically designed for the B2+Z1 platform: (1)~locomotion
        pre-training, (2)~manipulation pre-training with a stationary body, and
        (3)~joint fine-tuning for walk-while-grasp behaviour
        (\S\ref{sec:training_pipeline}).
  \item We demonstrate successful whole-body loco-manipulation of the B2+Z1
        system in Isaac Gym simulation, with the quadruped walking while
        simultaneously tracking arm end-effector targets (\S\ref{chp:exp}).
\end{enumerate}

\section{Thesis Structure}

The remainder of this report is organised as follows.
Chapter~\ref{chp:relwork} reviews related work on legged locomotion, arm
manipulation, and cooperative loco-manipulation.
Chapter~\ref{chp:method} describes the B2+Z1 platform, the adapted RoboDuet
framework, and the full RL training pipeline including policy architectures,
observation and action spaces, and reward design.
Chapter~\ref{chp:exp} presents experimental results from Isaac Gym simulation.
Chapter~\ref{chp:conclusion} concludes and discusses future work.
